% This file should be rename with the speaker's last name. (e.g : smith)

% Write the author names in the order you wish them to appear
% and underline the one who will give the talk.

% Only the speaker's email address is required.

\begin{abstract_online}{Proving and Disproving Subspaces with \emph{Mathematica}}{%
        \underline{Aharon Naiman}$^{1}$}{[naiman@jct.ac.il]
        }{%
        $^1$ Department of Applied Mathematics, Jerusalem College of
    Technology, Jerusalem, Israel}

Starting from the beginning of one's linear algebra education, one
ventures into the area of vector spaces.  After learning about the 10
axioms necessary for a vector space, the student delves into
subspaces.

As a subspace is also a vector space, we know that we can go back to
demonstrating that the 10 axioms are satisfied.  However, there is a
theorem that states that for a non-empty subset of a vector space, if
the subset is closed under vector addition and scalar multiplication,
then it is a subspace (and therefore also a vector space itself).

We present what tools are available to us in \emph{Mathematica} [1],
to assist in proving or disproving, in familiar mathematical notation,
whether a subset is a subspace.  First, we need to be able to model a
vector space, where the vector subset resides.  We demonstrate that we
cannot do this, for all trivial vector spaces studied in an elementary
course (e.g., function spaces).

Once we have the vector space, we present the necessary functions,
together with their many options, to assist in proving that a subset:
%
\begin{enumerate}
    \item is indeed a subspace---and \emph{why}, or alternatively,
    \item is \emph{not} a subspace, and use additional forms of the
          available functions to demonstrate \emph{intuitive}
          counterexamples.
\end{enumerate}

We develop proofs in both directions, using a few, different types of
vector spaces.

\textbf{Keywords} \newline{} linear algebra; education; proving and
disproving subspaces; \emph{Mathematica}

\textbf{References} \newline{}
\emph{Mathematica} at \texttt{www.wolfram.com/mathematica} \newline{}
\end{abstract_online}
    
